\documentclass[11pt,a4paper]{article}

\usepackage[margin=1.05in]{geometry}
\usepackage{amsmath,amssymb,amsthm,mathtools}
\usepackage{bm}
\usepackage[expansion=false]{microtype}
\usepackage{enumitem}
\usepackage{xcolor}
\usepackage[colorlinks=true,linkcolor=blue!55!black,citecolor=blue!55!black,
            urlcolor=blue!55!black]{hyperref}

%%%%%%%%%%%%%%%%%%%%%%%%%%%%%%%%%%%%%%%%%%%%%%%%%%%%%%%%%%%%%%%%%%%%%%%%%%%%%%%
%  Notation follows Bullo & Lewis, "Geometric Control of Mechanical Systems"
%  (TAM 49, Springer 2004) and Simpson-Porco & Bullo, "Contraction Theory on
%  Riemannian Manifolds".
%%%%%%%%%%%%%%%%%%%%%%%%%%%%%%%%%%%%%%%%%%%%%%%%%%%%%%%%%%%%%%%%%%%%%%%%%%%%%%%

\newcommand{\Gm}{\mathbb{G}}                      % kinetic energy metric
\newcommand{\Rm}{\mathbb{R}}
\newcommand{\ip}[2]{\langle\!\langle #1 , #2 \rangle\!\rangle}
\newcommand{\ipG}[2]{\langle\!\langle #1 , #2 \rangle\!\rangle_{\Gm}}
\newcommand{\dpair}[2]{\langle #1 ; #2 \rangle}   % canonical pairing
\newcommand{\Gsharp}{\Gm^{\sharp}}
\newcommand{\Gflat}{\Gm^{\flat}}
\newcommand{\Hess}{\operatorname{Hess}}
\newcommand{\grad}{\operatorname{grad}}
\newcommand{\dive}{\operatorname{div}}
\newcommand{\hor}{\operatorname{hor}}
\newcommand{\ver}{\operatorname{ver}}
\newcommand{\Sect}{\mathcal{K}}
\newcommand{\Jac}{\mathcal{R}}
\newcommand{\Damp}{\mathcal{D}}
\newcommand{\Stiff}{\Psi}
\newcommand{\Gsec}[1]{\Gamma^{\infty}(#1)}
\newcommand{\Cinf}[1]{C^{\infty}(#1)}
\newcommand{\id}{\mathrm{Id}}
\newcommand{\Lie}{\mathcal{L}}
\newcommand{\norm}[1]{\lVert #1 \rVert}
\newcommand{\normG}[1]{\lVert #1 \rVert_{\Gm}}
\newcommand{\Tstar}{T^{*}Q}
\newcommand{\Dt}{\tfrac{D}{dt}}

\DeclareMathOperator{\spec}{spec}
\DeclareMathOperator{\Vol}{Vol}
\DeclareMathOperator{\tr}{tr}

\theoremstyle{plain}
\newtheorem{theorem}{Theorem}[section]
\newtheorem{proposition}[theorem]{Proposition}
\newtheorem{lemma}[theorem]{Lemma}
\newtheorem{corollary}[theorem]{Corollary}
\theoremstyle{definition}
\newtheorem{definition}[theorem]{Definition}
\newtheorem{assumption}[theorem]{Assumption}
\theoremstyle{remark}
\newtheorem{remark}[theorem]{Remark}
\newtheorem{example}[theorem]{Example}

\title{\bfseries Contraction-Based Control of Simple Mechanical Systems:\\
A Cotangent-Bundle Formulation}
\author{Notes following Bullo--Lewis \cite{BL04} and Simpson-Porco--Bullo \cite{SB13}}
\date{}

\begin{document}
\maketitle

\begin{abstract}
\noindent
We set up the design of contraction-based feedback for a simple mechanical control
system $\Sigma = (Q,\Gm,V,\mathcal{F})$ in three stages. First we write the
uncontrolled dynamics on $\Tstar$ as a Hamiltonian vector field with respect to
the canonical symplectic form $\omega_0$ and verify that the drift is symplectic
(\S\ref{sec:cotangent}). Second we show that this is precisely the obstruction:
\emph{no} symplectic vector field can be contracting in the sense of
\cite[Def.~2.1]{SB13}, so any contraction-based design must deliberately destroy
$\omega_0$ (\S\ref{sec:obstruction}). Third we compute, in the splitting of
$T(\Tstar)$ induced by the Levi-Civita connection $\nabla$ of $\Gm$, the exact
covariant variational equations of the closed loop, show that block-diagonal
(Sasaki-type) contraction metrics fail \emph{identically} regardless of the
feedback, and give a sufficient condition for contraction using a metric with a
base--fiber cross term (\S\S\ref{sec:splitting}--\ref{sec:main}). The resulting
condition is
\[
  d \;>\; \frac{\sigma - \mu}{2\sqrt{\mu}},
\]
where $d$ is the damping gain and $[\mu,\sigma]$ brackets the spectrum of the
\emph{curvature-corrected shaped stiffness}
$\Stiff_\alpha = \Hess^{\sharp}(V+\varphi) + \Jac_v$, with $\Jac_v$ the Jacobi
operator of $(Q,\Gm)$. A short audit of which claims are quoted and which are
derived here appears in \S\ref{sec:audit}.
\end{abstract}

\tableofcontents

%%%%%%%%%%%%%%%%%%%%%%%%%%%%%%%%%%%%%%%%%%%%%%%%%%%%%%%%%%%%%%%%%%%%%%%%%%%%%%%
\section{Notation and standing setup}
\label{sec:notation}

We follow \cite[Ch.~3--4]{BL04} for the mechanics and \cite[\S1]{SB13} for the
contraction-theoretic objects.

\begin{itemize}[leftmargin=1.6em,itemsep=2pt]
\item $Q$ is a smooth, connected, $n$-dimensional \emph{configuration manifold};
      $\pi : \Tstar \to Q$ and $\tau_Q : TQ \to Q$ are the bundle projections.
\item $\Gm$ is the \emph{kinetic energy metric}, i.e.\ the Riemannian metric on
      $Q$ obtained from the inertia tensor \cite[\S4.2.3]{BL04}. We write
      $\ipG{\cdot}{\cdot}$ for the induced inner product on each $T_qQ$, and
      $\normG{\cdot}$ for the norm. The flat and sharp maps are
      $\Gflat : TQ \to \Tstar$, $\Gsharp = (\Gflat)^{-1}$; $\Gsharp$ also denotes
      the induced inner product $\ip{\cdot}{\cdot}_{\Gm^{-1}}$ on fibres of
      $\Tstar$.
\item $\nabla \equiv \overset{\Gm}{\nabla}$ is the Levi-Civita connection of
      $\Gm$; it is torsion-free and satisfies $\nabla\Gm = 0$, hence
      $\nabla\Gsharp = 0$ and $\nabla\Gflat = 0$.
\item $\dpair{\alpha_q}{v_q}$ is the canonical (metric-free) pairing between
      $T^{*}_qQ$ and $T_qQ$.
\item Curvature convention:
      $R(X,Y)Z = \nabla_X\nabla_Y Z - \nabla_Y\nabla_X Z - \nabla_{[X,Y]}Z$.
      Sectional curvature of the plane spanned by $u,v$:
      \[
        \Sect(u,v) \;=\;
        \frac{\ipG{R(u,v)v}{u}}{\normG{u}^2\normG{v}^2 - \ipG{u}{v}^2}.
      \]
\item For $\psi \in \Cinf{Q}$, $\grad\psi = \Gsharp\,d\psi$ and, as in
      \cite[\S5]{SB13},
      $\Hess\psi(v_q,w_q) = \ip{v_q}{\nabla_{w_q}\grad\psi}_{\Gm}$. We write
      $\Hess^{\sharp}\psi \coloneqq \nabla\grad\psi$ for the associated
      $\Gm$-symmetric $(1,1)$-tensor field, so that
      $\Hess\psi(u,w) = \ipG{\Hess^{\sharp}\psi(u)}{w}$.
\item $\Sigma = (Q,\Gm,V,\mathcal{F})$ is a \emph{simple mechanical control
      system} \cite[\S4.6.2]{BL04}: $V \in \Cinf{Q}$ is the potential and
      $\mathcal{F} = \{F^1,\dots,F^m\} \subset \Gsec{\Tstar}$ are the input
      covector fields (forces). The input vector fields are
      $Y_a = \Gsharp F^a$.
\end{itemize}

On the Lagrangian side, $\Sigma$ is the affine connection control system
\cite[\S4.6.3]{BL04}
\begin{equation}
  \nabla_{\gamma'(t)}\gamma'(t)
  \;=\; -\grad V(\gamma(t)) \;+\; \sum_{a=1}^{m} u^a(t)\, Y_a(\gamma(t)).
  \label{eq:accs}
\end{equation}
The object of these notes is a feedback rendering the closed loop
\emph{contracting} in the sense of \cite[Def.~2.1]{SB13}: there is a Riemannian
metric $G$ on the state manifold $M$ and $\lambda>0$ with
\begin{equation}
  \ip{\overset{G}{\nabla}_{w_x}X}{w_x}_{G(x)} \;\le\; -\lambda\,\norm{w_x}_{G(x)}^2
  \qquad \forall\, x \in \mathcal{U},\ \forall\, w_x \in T_xM .
  \label{eq:contraction}
\end{equation}

%%%%%%%%%%%%%%%%%%%%%%%%%%%%%%%%%%%%%%%%%%%%%%%%%%%%%%%%%%%%%%%%%%%%%%%%%%%%%%%
\section{The uncontrolled dynamics on \texorpdfstring{$\Tstar$}{T*Q} are symplectic}
\label{sec:cotangent}

\subsection{Legendre transform and Hamiltonian}

Because $\Gm$ is positive-definite, the Legendre transform
$\mathbb{F}L = \Gflat : TQ \to \Tstar$ is a fibrewise linear diffeomorphism
(hyperregularity). The Hamiltonian is
\begin{equation}
  H : \Tstar \to \Rm,
  \qquad
  H(\alpha_q) \;=\; \tfrac12\,\dpair{\alpha_q}{\Gsharp\alpha_q} + V(q)
  \;=\; \tfrac12\norm{\alpha_q}^2_{\Gm^{-1}} + V(q).
  \label{eq:H}
\end{equation}

\subsection{Canonical structures}

Let $\theta_0 \in \Gsec{T^{*}(\Tstar)}$ be the tautological (Liouville) one-form,
$\dpair{\theta_0(\alpha_q)}{w} = \dpair{\alpha_q}{T_{\alpha_q}\pi(w)}$, and set
$\omega_0 \coloneqq -d\theta_0$. In natural coordinates $(q^i,p_i)$ on $\Tstar$,
\begin{equation}
  \theta_0 = p_i\,dq^i,
  \qquad
  \omega_0 = dq^i \wedge dp_i .
  \label{eq:omega0}
\end{equation}
The Hamiltonian vector field $X_H \in \Gsec{T(\Tstar)}$ is defined by
$i_{X_H}\omega_0 = dH$, which in these coordinates is Hamilton's equations
$\dot q^i = \partial H/\partial p_i$, $\dot p_i = -\partial H/\partial q^i$.

\begin{proposition}[Symplecticity of the drift]
\label{prop:symplectic}
$\Lie_{X_H}\omega_0 = 0$. Consequently the flow $\Phi_t^{X_H}$ satisfies
$(\Phi_t^{X_H})^{*}\omega_0 = \omega_0$, and the Liouville volume
$\Lambda \coloneqq \tfrac{1}{n!}\,\omega_0^{\wedge n}$ is preserved:
$\dive_\Lambda X_H = 0$.
\end{proposition}

\begin{proof}
By Cartan's magic formula
$\Lie_{X_H}\omega_0 = d\,i_{X_H}\omega_0 + i_{X_H}\,d\omega_0$. The second term
vanishes since $d\omega_0 = -d^2\theta_0 = 0$; the first vanishes since
$i_{X_H}\omega_0 = dH$ is exact. Then
$\Lie_{X_H}\Lambda = \tfrac{1}{(n-1)!}(\Lie_{X_H}\omega_0)\wedge\omega_0^{\wedge(n-1)} = 0$.
\end{proof}

\subsection{Where the controls enter}

For $\beta \in T^{*}_qQ$ define the \emph{vertical lift}
$\ver_\alpha(\beta) = \frac{d}{d\varepsilon}\big|_{\varepsilon=0}(\alpha + \varepsilon\beta)
\in T_\alpha(\Tstar)$; in coordinates $\ver(\beta) = \beta_i\,\partial/\partial p_i$.
A direct computation from \eqref{eq:omega0} gives
\begin{equation}
  i_{\ver(\beta)}\omega_0 \;=\; -\pi^{*}\beta ,
  \label{eq:vertlift}
\end{equation}
so $\ver(\beta)$ is locally Hamiltonian if and only if $\beta$ is closed. The
controlled vector field on $\Tstar$ is
\begin{equation}
  X_\Sigma \;=\; X_H \;+\; \sum_{a=1}^{m} u^a\,\ver\!\left(F^a\right).
  \label{eq:Xsigma}
\end{equation}

\begin{remark}[The precise sense of ``the uncontrolled dynamics respect $\omega_0$'']
\label{rem:sense}
By Proposition~\ref{prop:symplectic}, $u \equiv 0$ gives a symplectic --- indeed
Hamiltonian --- flow. By \eqref{eq:vertlift}, an applied force $\beta$ keeps the
flow symplectic if and only if $\pi^{*}\beta$ is closed, i.e.\ iff the force is
locally a potential force \cite[\S4.4.4]{BL04}. \emph{Every} dissipative or
velocity-dependent feedback therefore breaks $\omega_0$. Section~\ref{sec:obstruction}
shows this is not an accident of the design but a hard requirement.
\end{remark}

%%%%%%%%%%%%%%%%%%%%%%%%%%%%%%%%%%%%%%%%%%%%%%%%%%%%%%%%%%%%%%%%%%%%%%%%%%%%%%%
\section{The obstruction: symplectic \texorpdfstring{$\Rightarrow$}{=>} not contracting}
\label{sec:obstruction}

The following makes precise why the design problem cannot be posed as a
structure-preserving one.

\begin{lemma}[Divergence bound]
\label{lem:div}
If $(\mathcal{U},X,G,\lambda)$ is a contracting system on a manifold $M$ of
dimension $N$, then $\dive_G X \le -N\lambda$ on $\mathcal{U}$, where $\dive_G$ is
the divergence with respect to the Riemannian volume of $G$.
\end{lemma}

\begin{proof}
Write $S \coloneqq \operatorname{sym}\big(\overset{G}{\nabla}X\big)$ for the
$G$-symmetric part of the $(1,1)$-tensor $\overset{G}{\nabla}X$. Condition
\eqref{eq:contraction} says exactly that $S \preceq -\lambda\,\id$. Taking traces
and using
\[
  \tr\big(\overset{G}{\nabla}X\big) = \tr(S) = \dive_G X
\]
gives the claim.
\end{proof}

\begin{proposition}[No symplectic vector field is contracting]
\label{prop:obstruction}
Let $(M,\omega)$ be a symplectic manifold of dimension $N = 2n \ge 2$, and let
$X \in \Gsec{TM}$ be symplectic, $\Lie_X\omega = 0$. Then there is no Riemannian
metric $G$, no $\lambda>0$ and no nonempty open $\mathcal{U} \subseteq M$ with
$\overline{\mathcal{U}}$ compact, forward $X$-invariant, and $X$ forward complete
on $\mathcal{U}$, such that $(\mathcal{U},X,G,\lambda)$ is contracting.
\end{proposition}

\begin{proof}
Write $\mu_G = f\,\Lambda$ with $\Lambda = \omega^{\wedge n}/n!$ and
$f \in \Cinf{M}$, $f>0$. Since $\Lie_X\Lambda = 0$ we have
$\dive_G X = \dive_\Lambda X + \Lie_X(\ln f) = \Lie_X(\ln f)$.
By Lemma~\ref{lem:div}, $\Lie_X(\ln f) \le -2n\lambda$ on $\mathcal{U}$, so along
any trajectory $t \mapsto \Phi_t(x)$ staying in $\mathcal{U}$,
\[
  \ln f(\Phi_t(x)) \;\le\; \ln f(x) - 2n\lambda t \;\xrightarrow[t\to\infty]{}\; -\infty.
\]
But $f$ attains a positive minimum on the compact set $\overline{\mathcal{U}}$,
so $\ln f$ is bounded below on $\mathcal{U}$ --- a contradiction.
\end{proof}

\begin{corollary}[Spectral form of the obstruction]
\label{cor:spectral}
Suppose in addition the hypotheses of \cite[Prop.~2.5(i)]{SB13} hold, so that $X$
has a fixed point $\bar x \in \mathcal{U}$. Evaluating \eqref{eq:contraction} at
$\bar x$, where $\overset{G}{\nabla}_{w}X = DX(\bar x)\cdot w$ is connection
independent, gives
$A^{\!\top}P + PA \preceq -2\lambda P$ with $A = DX(\bar x)$ and
$P = G(\bar x) \succ 0$; hence $A$ is Hurwitz. But a symplectic vector field is
locally Hamiltonian, so $A$ is a Hamiltonian matrix with respect to
$\omega(\bar x)$ and $\spec(A)$ is symmetric under $z \mapsto -z$. For $n \ge 1$
the two are incompatible.
\end{corollary}

\begin{remark}[Design consequence]
\label{rem:design}
Contraction is \emph{not} a structure-preserving property: it forces
$\Lie_{X_{\mathrm{cl}}}\omega_0 \ne 0$ and strict volume dissipation. The design
question is therefore not \emph{whether} to break $\omega_0$ but \emph{how}, and
--- equally important --- which metric $G$ on $\Tstar$ can certify the result.
Section~\ref{sec:sasaki} shows the obvious candidates cannot.
\end{remark}

%%%%%%%%%%%%%%%%%%%%%%%%%%%%%%%%%%%%%%%%%%%%%%%%%%%%%%%%%%%%%%%%%%%%%%%%%%%%%%%
\section{The \texorpdfstring{$\nabla$}{nabla}-splitting and the covariant variational equations}
\label{sec:splitting}

\subsection{Splitting of \texorpdfstring{$T(\Tstar)$}{T(T*Q)}}

Let $K_\alpha : T_\alpha(\Tstar) \to T^{*}_qQ$ be the connection map of $\nabla$:
if $s \mapsto \alpha(s)$ is a curve in $\Tstar$ over $s \mapsto q(s)$, then
$K(\alpha'(0)) = \tfrac{D}{ds}\alpha\big|_{0}$, the covariant derivative of the
covector field $\alpha$ along $q$. The pair
\begin{equation}
  T_\alpha(\Tstar) \;\ni\; w \;\longmapsto\;
  \big(u,\beta\big) \coloneqq \big(T_\alpha\pi(w),\, K_\alpha(w)\big)
  \;\in\; T_qQ \times T^{*}_qQ
  \label{eq:splitting}
\end{equation}
is a linear isomorphism, splitting $T_\alpha(\Tstar) = H_\alpha \oplus V_\alpha$
into horizontal and vertical subspaces. Throughout we also write
$\xi \coloneqq \Gsharp\beta \in T_qQ$, so that
$\dpair{\beta}{u} = \ipG{\xi}{u}$.

\begin{lemma}[Canonical objects in the splitting]
\label{lem:splitting}
With $\nabla$ torsion-free and $w_j \leftrightarrow (u_j,\beta_j)$,
\begin{align}
  \theta_0(w) &= \dpair{\alpha}{u}, \label{eq:theta-split}\\
  \omega_0(w_1,w_2) &= \dpair{\beta_2}{u_1} - \dpair{\beta_1}{u_2},
     \label{eq:omega-split}\\
  X_H(\alpha) &= \hor_\alpha\!\big(\Gsharp\alpha\big) \;-\; \ver_\alpha\!\big(dV(q)\big).
     \label{eq:XH-split}
\end{align}
Equivalently, with $v \coloneqq \Gsharp\alpha$, the uncontrolled dynamics read
\begin{equation}
  \dot q = v = \Gsharp\alpha,
  \qquad
  \Dt \alpha = -dV(q),
  \label{eq:ham-covariant}
\end{equation}
and the controlled dynamics \eqref{eq:Xsigma} read
$\Dt\alpha = -dV(q) + \sum_a u^a F^a(q)$.
\end{lemma}

\begin{proof}[Proof sketch]
\eqref{eq:theta-split} is the definition of $\theta_0$. For
\eqref{eq:omega-split}, evaluate at a point in normal coordinates for $\nabla$
(all $\Gamma^k_{ij}$ vanish there), where the splitting \eqref{eq:splitting}
coincides with the coordinate splitting; then
$(dq^i\wedge dp_i)(w_1,w_2) = u_1^i(\beta_2)_i - u_2^i(\beta_1)_i$. Torsion-freeness
is what guarantees no extra term appears at points where $\Gamma \ne 0$.
For \eqref{eq:ham-covariant}: the horizontal component is
$\dot q^i = \Gm^{ij}p_j$ by \eqref{eq:H}, and the vertical component is
\[
  \big(\Dt\alpha\big)_i
  = \dot p_i - \Gamma^k_{ij}\dot q^{\,j} p_k
  = -\frac{\partial V}{\partial q^i}
    - \tfrac12\frac{\partial \Gm^{jk}}{\partial q^i}p_jp_k
    - \Gamma^k_{ij}\Gm^{j\ell}p_\ell p_k .
\]
Metric compatibility $\nabla\Gm^{-1}=0$ gives
$\partial_i\Gm^{jk} = -\Gamma^j_{i\ell}\Gm^{\ell k} - \Gamma^k_{i\ell}\Gm^{j\ell}$,
and substituting cancels the last two terms identically, leaving
$\Dt\alpha = -dV$. Applying $\Gsharp$ (which is $\nabla$-parallel) recovers
\eqref{eq:accs} with $u\equiv0$.
\end{proof}

\begin{remark}[The symmetric partner of $\omega_0$]
\label{rem:partner}
By \eqref{eq:omega-split}, the canonical pairing appears \emph{antisymmetrized}
in $\omega_0$. Its \emph{symmetrization}
\begin{equation}
  \Pi(w_1,w_2) \coloneqq \dpair{\beta_2}{u_1} + \dpair{\beta_1}{u_2}
  \label{eq:Pi}
\end{equation}
is a metric-independent symmetric bilinear form on $T_\alpha(\Tstar)$ of signature
$(n,n)$. It is exactly the ingredient the contraction metric must contain
(Theorem~\ref{thm:main}); Proposition~\ref{prop:sasaki} shows a metric without it
cannot work.
\end{remark}

\subsection{The variational (Jacobi) system}

Let $(s,t)\mapsto \alpha(s,t)$ be a smooth family of trajectories of the closed
loop, with base curves $q(s,t)$ and $v = \partial_t q = \Gsharp\alpha$. Set
\begin{equation}
  u \coloneqq \partial_s q,
  \qquad
  \beta \coloneqq \tfrac{D}{ds}\alpha,
  \qquad
  \xi \coloneqq \Gsharp\beta .
\end{equation}
Then $(u,\beta)$ is precisely the splitting \eqref{eq:splitting} of the variation
field $S(s,t) = \tfrac{d}{ds}\Phi_t(\gamma(s))$ appearing in the proof of
\cite[Thm.~2.3]{SB13}.

\begin{proposition}[Covariant variational equations]
\label{prop:variational}
For the controlled dynamics $\Dt\alpha = -dV + \sum_a u^aF^a$,
\begin{align}
  \Dt u &= \xi, \label{eq:var1}\\
  \Dt \xi &= -\Hess^{\sharp}V(u) \;-\; \Jac_v(u) \;+\; \delta_u,
    \label{eq:var2}
\end{align}
where $\Jac_v \in \Gsec{T^1_1Q}$ is the \emph{Jacobi} (tidal) operator
\begin{equation}
  \Jac_v(u) \coloneqq R(u,v)v ,
  \qquad
  \ipG{\Jac_v u}{u} = \Sect(u,v)\left(\normG{u}^2\normG{v}^2 - \ipG{u}{v}^2\right),
  \label{eq:jacobi-op}
\end{equation}
which is $\Gm$-symmetric, and
$\delta_u = \Gsharp\tfrac{D}{ds}\!\left(\sum_a u^aF^a\right)$ is the variation of
the control force.
\end{proposition}

\begin{proof}
\eqref{eq:var1}: since $\nabla$ is torsion-free,
$\tfrac{D}{ds}\partial_t q = \tfrac{D}{dt}\partial_s q$, and since
$\nabla\Gsharp = 0$,
$\tfrac{D}{ds}(\Gsharp\alpha) = \Gsharp\tfrac{D}{ds}\alpha = \Gsharp\beta = \xi$.
Hence $\Dt u = \xi$.

\eqref{eq:var2}: for any vector field $W$ along the map $(s,t)\mapsto q(s,t)$,
differentiating the function $\dpair{\alpha}{W}$ twice and using
$(\tfrac{D}{ds}\tfrac{D}{dt} - \tfrac{D}{dt}\tfrac{D}{ds})W = R(u,v)W$ gives
\[
  \dpair{\left(\tfrac{D}{ds}\tfrac{D}{dt} - \tfrac{D}{dt}\tfrac{D}{ds}\right)\alpha}{W}
  \;=\; -\dpair{\alpha}{R(u,v)W}.
\]
Therefore
$\Dt\beta = \tfrac{D}{ds}\!\left(\Dt\alpha\right) + \alpha\circ R(u,v)$.
Now $\tfrac{D}{ds}(-dV) = -\nabla_u dV = -\Gflat\big(\Hess^{\sharp}V(u)\big)$, and
for any $W$,
\[
  \dpair{\alpha\circ R(u,v)}{W} = \ipG{R(u,v)W}{v} = -\ipG{R(u,v)v}{W},
\]
so $\Gsharp\big(\alpha\circ R(u,v)\big) = -R(u,v)v = -\Jac_v(u)$. Applying
$\Gsharp$ throughout yields \eqref{eq:var2}. The expression \eqref{eq:jacobi-op}
for $\ipG{\Jac_vu}{u}$ is the definition of sectional curvature; $\Gm$-symmetry of
$\Jac_v$ follows from the pair symmetry $\ipG{R(X,Y)Z}{W} = \ipG{R(Z,W)X}{Y}$.
\end{proof}

\begin{remark}[Consistency check]
With $V \equiv 0$ and $\delta_u \equiv 0$, \eqref{eq:var1}--\eqref{eq:var2} become
$\tfrac{D^2u}{dt^2} + R(u,v)v = 0$, the Jacobi equation of geodesic deviation, as
it must: the uncontrolled unforced dynamics of $\Sigma$ are the geodesic spray of
$\Gm$.
\end{remark}

\begin{remark}[Why this suffices for \eqref{eq:contraction}]
\label{rem:equivalence}
Equation~(8) in the proof of \cite[Thm.~2.3]{SB13} establishes
$\tfrac{d}{dt}\norm{S(s,t)}_G^2 = 2\ip{\overset{G}{\nabla}_{S}X}{S}_G$.
Hence \eqref{eq:contraction} holds on $\mathcal{U}$ if and only if
$\tfrac{d}{dt}\norm{w}_G^2 \le -2\lambda\norm{w}_G^2$ for every variation field
$w$ along every trajectory in $\mathcal{U}$. This is what we verify below, and it
spares us computing the Levi-Civita connection of $G$ on the $2n$-manifold
$\Tstar$ explicitly.
\end{remark}

%%%%%%%%%%%%%%%%%%%%%%%%%%%%%%%%%%%%%%%%%%%%%%%%%%%%%%%%%%%%%%%%%%%%%%%%%%%%%%%
\section{Block-diagonal metrics cannot certify contraction}
\label{sec:sasaki}

The first candidate metric on $\Tstar$ is the Sasaki-type lift: horizontal block
$\Gm$, vertical block $\Gm^{-1}$, no cross term. More generally consider any
metric that is block diagonal and fibre-independent in the splitting
\eqref{eq:splitting}.

\begin{proposition}[Obstruction to block-diagonal certificates]
\label{prop:sasaki}
Let $A,C \in \Gsec{T^1_1Q}$ be $\Gm$-symmetric and positive-definite, and define
the Riemannian metric on $\Tstar$
\[
  G_\alpha(w,w) = \ipG{A(q)u}{u} + \ipG{C(q)\xi}{\xi},
  \qquad w \leftrightarrow (u,\xi).
\]
Suppose $\mathcal{W}\subseteq\Tstar$ contains a pair $\pm\alpha_q$ with
$\alpha_q \ne 0$. Then $G$ cannot satisfy \eqref{eq:contraction} on $\mathcal{W}$
for the closed loop, \emph{for any feedback whatsoever}.
\end{proposition}

\begin{proof}
Differentiating along a trajectory with velocity $v$, using
$\tfrac{d}{dt}\ipG{Au}{u} = \ipG{(\nabla_vA)u}{u} + 2\ipG{A\,\Dt u}{u}$,
\[
  \tfrac12\tfrac{d}{dt}G(w,w)
  = \tfrac12\ipG{(\nabla_vA)u}{u} + \ipG{A\xi}{u}
    + \tfrac12\ipG{(\nabla_vC)\xi}{\xi} + \ipG{C\,\Dt\xi}{\xi},
\]
where we used $\Dt u = \xi$ from \eqref{eq:var1}. Since $(u,\xi)$ ranges over all
of $T_qQ\times T_qQ$, evaluate at a variation with $\xi = 0$ and $u \ne 0$: all
terms containing $\xi$ drop, \emph{including the entire dependence on
$\Dt\xi$ and therefore on the feedback}, leaving
\[
  \tfrac12\tfrac{d}{dt}G(w,w)\Big|_{\xi=0} = \tfrac12\ipG{(\nabla_vA)u}{u}.
\]
Condition \eqref{eq:contraction} would require
$\ipG{(\nabla_vA)u}{u} \le -2\lambda\ipG{Au}{u}$. Applying this at both
$\alpha_q$ and $-\alpha_q$ (i.e.\ at $v$ and $-v$) and using linearity of
$\nabla_vA$ in $v$ gives simultaneously
$\ipG{(\nabla_vA)u}{u} \le -2\lambda\ipG{Au}{u}$ and
$\ipG{(\nabla_vA)u}{u} \ge 2\lambda\ipG{Au}{u}$, whence
$\ipG{Au}{u}\le 0$ --- contradicting $A \succ 0$, $u\ne0$, $\lambda>0$.
\end{proof}

\begin{corollary}
The Sasaki metric ($A = \id$, $C = \id$, both $\nabla$-parallel) never certifies
contraction. Neither does the total-energy metric on $\Tstar$.
\end{corollary}

This is the intrinsic counterpart of the familiar linear fact: for
$M\ddot q + D\dot q + Kq = 0$ with $P = \operatorname{diag}(K, M^{-1})$ (the energy
metric), $A^{\!\top}P + PA = \operatorname{diag}(0, -2M^{-1}DM^{-1})$ is only
negative \emph{semi}definite. It is also exactly why Example~1 of \cite{SB13}
carries the cross term $b\varepsilon\,dx\otimes dy$.

%%%%%%%%%%%%%%%%%%%%%%%%%%%%%%%%%%%%%%%%%%%%%%%%%%%%%%%%%%%%%%%%%%%%%%%%%%%%%%%
\section{Contraction metrics with a base--fibre cross term}
\label{sec:main}

\subsection{The metric}

\begin{definition}[Skewed lift]
\label{def:metric}
Let $a,b,c \in \Rm$ with $a,c>0$ and $ac > b^2$. Define $G$ on $\Tstar$ by
\begin{equation}
  G_\alpha(w,w) \;=\; a\,\normG{u}^2 \;+\; 2b\,\ipG{u}{\xi} \;+\; c\,\normG{\xi}^2,
  \qquad w \leftrightarrow (u,\xi = \Gsharp\beta).
  \label{eq:G}
\end{equation}
\end{definition}

$G$ is a genuine Riemannian metric on the $2n$-manifold $\Tstar$: the horizontal
block is $a\,\pi^{*}\Gm$, the vertical block is $c\,\Gm^{-1}$, and the cross term
is $b\,\Pi$ with $\Pi$ the canonical symmetric form \eqref{eq:Pi} of
Remark~\ref{rem:partner}. Metrics of this shape belong to the class called
\emph{$g$-natural} in the Riemannian-geometry literature (here with constant,
fibre-independent coefficients).

\subsection{The feedback}

\begin{assumption}[Full actuation]
\label{ass:actuation}
$\operatorname{span}_{\Rm}\{F^1(q),\dots,F^m(q)\} = T^{*}_qQ$ for each $q$ in the
region of interest.
\end{assumption}

Under Assumption~\ref{ass:actuation} we may realise any desired control force. Take
the potential-shaping-plus-damping law of \cite[\S10.4, \S11.1]{BL04}:
\begin{equation}
  F_{\mathrm{ctrl}}(\alpha) \;=\; -\,d\varphi(q) \;-\; \Gflat\big(\Damp(q)\,\Gsharp\alpha\big),
  \label{eq:feedback}
\end{equation}
with $\varphi \in \Cinf{Q}$ the shaping potential and $\Damp \in \Gsec{T^1_1Q}$
$\Gm$-symmetric and positive-definite. The closed loop is
\begin{equation}
  \dot q = v = \Gsharp\alpha,
  \qquad
  \Dt\alpha = -\,d(V+\varphi)(q) - \Gflat(\Damp v),
  \qquad\text{i.e.}\qquad
  \nabla_v v = -\grad(V+\varphi) - \Damp v.
  \label{eq:closedloop}
\end{equation}
By Remark~\ref{rem:sense} this is not symplectic --- as
Proposition~\ref{prop:obstruction} demands.

\begin{assumption}[Parallel damping]
\label{ass:damping}
$\nabla\Damp = 0$. This holds in particular for $\Damp = d\,\id$ with
$d \in \Rm_{>0}$.
\end{assumption}

Under Assumptions~\ref{ass:actuation}--\ref{ass:damping}, substituting
\eqref{eq:closedloop} into Proposition~\ref{prop:variational} gives
$\delta_u = -\Hess^{\sharp}\varphi(u) - \Damp\xi$, hence the closed-loop
variational system
\begin{equation}
  \boxed{\quad
  \Dt u = \xi,
  \qquad
  \Dt \xi = -\,\Stiff_\alpha(u) \;-\; \Damp\,\xi,
  \qquad
  \Stiff_\alpha \coloneqq \Hess^{\sharp}(V+\varphi)(q) + \Jac_v .
  \quad}
  \label{eq:cl-variational}
\end{equation}
$\Stiff_\alpha$ is $\Gm$-symmetric; we call it the \emph{curvature-corrected
shaped stiffness}. Note it depends on the full state $\alpha$, not just on $q$.

\subsection{Main result}

\begin{theorem}[Contraction certificate]
\label{thm:main}
Let Assumptions~\ref{ass:actuation}--\ref{ass:damping} hold with
$\Damp = d\,\id$, $d>0$. Let $\mathcal{W}\subseteq\Tstar$ be open, forward
invariant for \eqref{eq:closedloop}, $K$-reachable in the sense of
\cite[Def.~2.2]{SB13}, with the closed-loop field forward complete on
$\mathcal{W}$. Suppose there exist $0 < \mu \le \sigma$ with
\begin{equation}
  \mu\,\normG{u}^2 \;\le\; \ipG{\Stiff_\alpha u}{u}
  \quad\text{and}\quad
  \normG{\Stiff_\alpha u} \;\le\; \sigma\,\normG{u}
  \qquad \forall\,\alpha\in\mathcal{W},\ \forall\,u\in T_qQ,
  \label{eq:mu-sigma}
\end{equation}
i.e.\ $\spec(\Stiff_\alpha)\subset[\mu,\sigma]$ uniformly on $\mathcal{W}$. If
\begin{equation}
  \boxed{\;\; d \;>\; \frac{\sigma-\mu}{2\sqrt{\mu}} \;\;}
  \label{eq:main-condition}
\end{equation}
then $G$ of Definition~\ref{def:metric} with
\begin{equation}
  c = 1, \qquad b = \tfrac{d}{2}, \qquad a = \tfrac{d^2}{2} + \tfrac{\mu+\sigma}{2}
  \label{eq:gains}
\end{equation}
makes $(\mathcal{W}, X_{\mathrm{cl}}, G, \lambda)$ a contracting system, with
\begin{equation}
  \lambda \;=\; \frac{\lambda_{\min}(\mathcal{Q})}{\lambda_{\max}(\mathcal{P})},
  \qquad
  \mathcal{Q} = \begin{bmatrix} \tfrac{\mu d}{2} & -\tfrac{\kappa}{2}\\[2pt]
                                -\tfrac{\kappa}{2} & \tfrac{d}{2}\end{bmatrix},
  \quad
  \mathcal{P} = \begin{bmatrix} a & b\\ b & c\end{bmatrix},
  \quad
  \kappa \le \tfrac{\sigma-\mu}{2},
  \label{eq:rate}
\end{equation}
and consequently, by \cite[Thm.~2.3]{SB13},
\[
  d_G\big(\Phi_t(\alpha_0),\Phi_t(\alpha_1)\big) \;\le\;
  K e^{-\lambda t}\, d_G(\alpha_0,\alpha_1),
  \qquad \forall\,\alpha_0,\alpha_1\in\mathcal{W},\ t\ge0 .
\]
\end{theorem}

\begin{proof}
By Remark~\ref{rem:equivalence} it suffices to bound
$\tfrac{d}{dt}G(w,w)$ along \eqref{eq:cl-variational}. Since $a,b,c$ are constants
and $\Gm$ is $\nabla$-parallel,
\begin{align*}
  \tfrac12\tfrac{d}{dt}G(w,w)
   &= a\ipG{\Dt u}{u} + b\big(\ipG{\Dt u}{\xi} + \ipG{u}{\Dt \xi}\big)
      + c\ipG{\Dt \xi}{\xi}\\
   &= a\ipG{\xi}{u} + b\normG{\xi}^2
      + b\ipG{u}{-\Stiff u - \Damp\xi}
      + c\ipG{-\Stiff u - \Damp\xi}{\xi}\\
   &= -\,b\ipG{\Stiff u}{u} \;-\; c\ipG{\Damp\xi}{\xi} \;+\; b\normG{\xi}^2
      \;+\; \ipG{\big(a\,\id - b\Damp - c\Stiff\big)\xi}{u},
\end{align*}
where $\Gm$-symmetry of $\Damp$ and $\Stiff$ was used to collect the cross terms.
With $\Damp = d\,\id$ and $\kappa \coloneqq
\sup_{\mathcal{W}}\norm{a\,\id - bd\,\id - c\,\Stiff_\alpha}$, the bounds
\eqref{eq:mu-sigma} give
\[
  \tfrac12\tfrac{d}{dt}G(w,w)
  \;\le\; -\,b\mu\normG{u}^2 - (cd-b)\normG{\xi}^2 + \kappa\normG{u}\normG{\xi}
  \;=\; -\,z^{\!\top}\mathcal{Q}\,z,
  \qquad z = \begin{bmatrix}\normG{u}\\ \normG{\xi}\end{bmatrix},
\]
with $\mathcal{Q} = \begin{bsmallmatrix} b\mu & -\kappa/2\\ -\kappa/2 & cd-b\end{bsmallmatrix}$.
Choosing $c=1$ and $a = bd + \tfrac{\mu+\sigma}{2}$ centres the residual operator,
$a\,\id - bd\,\id - \Stiff = \tfrac{\mu+\sigma}{2}\id - \Stiff$, so
$\kappa \le \tfrac{\sigma-\mu}{2}$ by \eqref{eq:mu-sigma}. Then
$\mathcal{Q}\succ0$ iff $b < d$ and $b\mu(d-b) > \kappa^2/4$. The left side is
maximised at $b = d/2$, giving $\mu d^2/4 > (\sigma-\mu)^2/16$, i.e.\
\eqref{eq:main-condition}. With $b=d/2$ we get
$a = d^2/2 + (\mu+\sigma)/2 > d^2/4 = b^2 = b^2/c$, so $G$ is positive-definite as
required in Definition~\ref{def:metric}. Finally
$G(w,w) = z^{\!\top}\mathcal{P}z \le \lambda_{\max}(\mathcal{P})\norm{z}^2$ and
$z^{\!\top}\mathcal{Q}z \ge \lambda_{\min}(\mathcal{Q})\norm{z}^2$, so
$\tfrac{d}{dt}G(w,w) \le -2\lambda\,G(w,w)$ with $\lambda$ as in \eqref{eq:rate}.
The stated bound is then \cite[Thm.~2.3]{SB13}.
\end{proof}

\begin{remark}[Isotropic case: any damping works]
\label{rem:isotropic}
If $\Stiff_\alpha = \mu\,\id$ on $\mathcal{W}$ then $\sigma = \mu$,
$\kappa = 0$, and \eqref{eq:main-condition} reads $d > 0$: \emph{arbitrarily small
damping suffices}, provided the metric carries the cross term $b = d/2$. This is
the intrinsic version of Example~1 of \cite{SB13}, where the damped oscillator is
shown contracting for every $b>0$ with a metric containing
$b\varepsilon\,dx\otimes dy$. (The coefficients differ from theirs --- contraction
metrics are very far from unique --- but the essential feature, a nonzero cross
term, is the same.)
\end{remark}

\begin{remark}[Non-parallel damping]
Dropping Assumption~\ref{ass:damping}, \eqref{eq:cl-variational} acquires the
extra term $\mathcal{N}_v(u) \coloneqq (\nabla_u\Damp)v$, which is not
$\Gm$-symmetric. Splitting $\mathcal{N}_v$ into symmetric and skew parts, the
symmetric part is absorbed into $\Stiff_\alpha$ (revising $\mu,\sigma$) and the
skew part contributes only to the cross term, i.e.\ to $\kappa$. The structure of
Theorem~\ref{thm:main} is unchanged; only the constants degrade.
\end{remark}

\begin{remark}[State-dependent coefficients]
Allowing $a,b,c$ in \eqref{eq:G} to be functions on $Q$, or $\Gm$-symmetric tensor
fields $A,B,C$, adds $\tfrac12\ipG{(\nabla_vA)u}{u}$,
$\ipG{(\nabla_vB)u}{\xi}$, $\tfrac12\ipG{(\nabla_vC)\xi}{\xi}$ to the derivative.
Note these are derivatives along the \emph{trajectory} velocity $v$, not along the
variation --- a consequence of doing the whole computation covariantly. This extra
freedom is what one would need to handle strongly anisotropic $\Stiff_\alpha$
(where the constant-coefficient bound $\kappa \le (\sigma-\mu)/2$ is
conservative), at the cost of a coupled PDE-like condition on $A,B,C$.
\end{remark}

%%%%%%%%%%%%%%%%%%%%%%%%%%%%%%%%%%%%%%%%%%%%%%%%%%%%%%%%%%%%%%%%%%%%%%%%%%%%%%%
\section{Reading the condition: curvature, velocity, and the size of \texorpdfstring{$\mathcal{W}$}{W}}
\label{sec:reading}

By \eqref{eq:cl-variational} and \eqref{eq:jacobi-op},
\begin{equation}
  \ipG{\Stiff_\alpha u}{u}
  \;=\; \Hess(V+\varphi)(u,u)
  \;+\; \Sect(u,v)\left(\normG{u}^2\normG{v}^2 - \ipG{u}{v}^2\right).
  \label{eq:stiffness-decomp}
\end{equation}
Three consequences.

\paragraph{Positive curvature helps.} Where $\Sect \ge 0$, the curvature term
\emph{adds} to the effective stiffness. This is the geodesic-deviation picture:
on a positively curved manifold neighbouring geodesics focus, and the focusing
does part of the controller's work. On a Lie group with a bi-invariant metric,
$\Sect(X,Y) = \tfrac14\normG{[X,Y]}^2 \ge 0$ for $\Gm$-orthonormal $X,Y$; e.g.\
a rigid body with inertia tensor $\mathbb{J} = j\,\id$ on $SO(3)$.

\paragraph{Negative curvature bounds $\mathcal{W}$ in velocity.} If
$\Sect \ge \Sect_{-}$ with $\Sect_{-} < 0$, then from
\eqref{eq:stiffness-decomp}
\[
  \mu \;\ge\; \lambda_{\min}\!\big(\Hess^{\sharp}(V+\varphi)\big)
              \;-\; |\Sect_{-}|\,\sup_{\mathcal{W}}\normG{v}^2 .
\]
Positivity of $\mu$ --- which \eqref{eq:main-condition} requires --- therefore
forces $\sup_{\mathcal{W}}\normG{v}^2 <
\lambda_{\min}(\Hess^{\sharp}(V+\varphi))/|\Sect_{-}|$. Crucially,
\emph{damping cannot repair this}: $\Damp$ enters only the $\xi$-block of
\eqref{eq:cl-variational}, whereas $\mu$ lives in the $u$-block. On a negatively
curved configuration manifold, contraction regions are intrinsically
velocity-bounded, and enlarging them requires stiffer potential shaping, not more
damping. For a rigid body on $SO(3)$ with an asymmetric inertia tensor the metric
$\Gm$ is left- but not bi-invariant, and sectional curvatures of some planes are
negative --- so this case is live, not academic.

\paragraph{Topology bounds $\mathcal{W}$ too.} \cite[\S2]{SB13} notes, following
Bhat--Bernstein, that a contraction region is contractible. Since $\Tstar$
deformation-retracts onto $Q$, $\mathcal{W}$ can never be all of $\Tstar$ when $Q$
is not contractible ($Q = S^2$, $SO(3)$, $\mathbb{T}^n$, \dots). Contraction-based
designs on such $Q$ are unavoidably local, in the same way and for the same reason
that smooth global stabilisation is \cite[\S10.5.7]{BL04}.

%%%%%%%%%%%%%%%%%%%%%%%%%%%%%%%%%%%%%%%%%%%%%%%%%%%%%%%%%%%%%%%%%%%%%%%%%%%%%%%
\section{Design recipe}
\label{sec:recipe}

\begin{enumerate}[leftmargin=1.9em,itemsep=3pt]
\item Model $\Sigma = (Q,\Gm,V,\mathcal{F})$ as in \cite[Ch.~4]{BL04}; compute
      $\nabla$ and $R$ for $\Gm$.
\item Write the drift on $\Tstar$ via \eqref{eq:H}--\eqref{eq:omega0}; confirm
      $\Lie_{X_H}\omega_0 = 0$ (Proposition~\ref{prop:symplectic}). This is the
      structure the feedback must break.
\item Verify Assumption~\ref{ass:actuation}. If it fails, this construction does
      not apply --- see \S\ref{sec:limits}.
\item Pick a shaping potential $\varphi$ with $\Hess^{\sharp}(V+\varphi) \succ 0$
      on the region of interest and a unique critical point at the target
      \cite[\S10.4, \S11.1.1]{BL04}.
\item Choose a candidate velocity bound $\bar v$; compute
      $\mu,\sigma$ for $\Stiff_\alpha = \Hess^{\sharp}(V+\varphi) + \Jac_v$ over
      $\{\normG{v}\le\bar v\}$ using \eqref{eq:stiffness-decomp}. If $\mu \le 0$,
      shrink $\bar v$ or stiffen $\varphi$.
\item Choose $d > (\sigma-\mu)/(2\sqrt{\mu})$ and set
      $(a,b,c)$ by \eqref{eq:gains}. The certificate is $G$ of
      \eqref{eq:G}; the rate is \eqref{eq:rate}.
\item Verify forward invariance and $K$-reachability of $\mathcal{W}$, e.g.\ by
      taking $\mathcal{W}$ a sublevel set of the closed-loop energy
      $\tfrac12\normG{v}^2 + (V+\varphi)$ contained in
      $\{\normG{v}\le\bar v\}$; sublevel sets of a strict Lyapunov function are
      forward invariant, and geodesic balls are $1$-reachable
      \cite[Prop.~2.5(iii)]{SB13}.
\end{enumerate}

%%%%%%%%%%%%%%%%%%%%%%%%%%%%%%%%%%%%%%%%%%%%%%%%%%%%%%%%%%%%%%%%%%%%%%%%%%%%%%%
\section{Limitations and open directions}
\label{sec:limits}

\begin{itemize}[leftmargin=1.6em,itemsep=3pt]
\item \textbf{Full actuation is used essentially.} Under-actuated $\Sigma$
      cannot realise \eqref{eq:feedback}. The relevant structural theory is then
      kinematic reduction and decoupling vector fields \cite[Ch.~8]{BL04}, and
      contraction certificates in that setting are, as far as these notes go,
      open. Note also that under-actuation is exactly where the symplectic
      structure is hardest to break usefully: the reachable directions in which
      $\omega_0$ can be destroyed are constrained by
      $\operatorname{span}\{F^a\}$.
\item \textbf{The conclusion of \cite{SB13} is not resolved.} That paper closes by
      flagging as open ``finding succinct and meaningful conditions for a simple
      mechanical system to exhibit contracting behavior.'' Theorem~\ref{thm:main}
      gives one such condition \emph{for the closed loop under full actuation};
      it says nothing about when an uncontrolled or under-actuated simple
      mechanical system is contracting. Indeed
      Proposition~\ref{prop:obstruction} shows that the uncontrolled system never
      is.
\item \textbf{Conservatism.} The constant-coefficient bound
      $\kappa \le (\sigma-\mu)/2$ is tight only when $\Stiff_\alpha$ is close to
      isotropic. For strongly anisotropic inertia/curvature, tensor-valued
      $A,B,C$ should be expected to give much larger $\mathcal{W}$.
\item \textbf{Cut locus.} $d_G$ on $\Tstar$ is a genuine distance, but the
      exponential bound is stated in $d_G$, not in any coordinate norm; converting
      to a usable estimate requires control of the cut locus of $G$, which for a
      $g$-natural metric is not determined by that of $\Gm$.
\item \textbf{Alternative: keep the structure, change the notion.} If preserving
      $\omega_0$ matters (variational integrators, energy methods), contraction is
      the wrong target by Proposition~\ref{prop:obstruction}; one would instead
      use relative equilibria and energy--momentum methods
      \cite[\S6.3]{BL04}, or contraction only in a quotient / on a
      symplectically reduced space.
\end{itemize}

%%%%%%%%%%%%%%%%%%%%%%%%%%%%%%%%%%%%%%%%%%%%%%%%%%%%%%%%%%%%%%%%%%%%%%%%%%%%%%%
\section{Audit of claims}
\label{sec:audit}

Because it matters which statements here are quoted and which are constructed:

\begin{itemize}[leftmargin=1.6em,itemsep=3pt]
\item \emph{Quoted verbatim in substance from \cite{SB13}:} Definition 2.1
      (contracting system), Definition 2.2 ($K$-reachability), Theorem 2.3 and
      its equation (8), Proposition 2.5 (existence of a fixed point, Krasovski\u{\i}
      function, incremental Lyapunov function, volume contraction), Example 1
      (damped oscillator with cross-term metric), the remark that contraction
      regions are contractible, and the open problem quoted in
      \S\ref{sec:limits}.
\item \emph{Quoted in substance from \cite{BL04}:} the affine connection control
      system \eqref{eq:accs} (\S4.6.3), the kinetic-energy-as-metric construction
      (\S4.2.3), potential forces (\S4.4.4), PD and potential shaping
      (\S\S10.4, 11.1), kinematic reduction (Ch.~8), obstructions to smooth
      stabilisation (\S10.5.7).
\item \emph{Standard textbook differential geometry, stated without proof:}
      Cartan's magic formula; the tautological one-form and $\omega_0 = -d\theta_0$;
      the symmetries of $R$; the Jacobi equation; nonnegative sectional curvature
      of bi-invariant metrics.
\item \emph{Derived here and worth independent checking:} Lemma~\ref{lem:splitting}
      \eqref{eq:omega-split}--\eqref{eq:ham-covariant};
      Proposition~\ref{prop:variational} (in particular the sign of the curvature
      term); Proposition~\ref{prop:obstruction} and Corollary~\ref{cor:spectral};
      Proposition~\ref{prop:sasaki}; Theorem~\ref{thm:main} and the constant
      \eqref{eq:main-condition}. The consistency checks available are: the
      $V\equiv0$ reduction to the Jacobi equation, and the $\sigma=\mu$ reduction
      of \eqref{eq:main-condition} to ``any $d>0$'', matching Example~1 of
      \cite{SB13}.
\item \emph{Attributed loosely:} the term ``$g$-natural metric'' is used in the
      Riemannian-geometry sense (Kowalski--Sekizawa and successors); the precise
      classification of that class is not used here, only the shape \eqref{eq:G}.
\end{itemize}

%%%%%%%%%%%%%%%%%%%%%%%%%%%%%%%%%%%%%%%%%%%%%%%%%%%%%%%%%%%%%%%%%%%%%%%%%%%%%%%
\begin{thebibliography}{9}

\bibitem{BL04}
F.~Bullo and A.~D.~Lewis,
\emph{Geometric Control of Mechanical Systems: Modeling, Analysis, and Design for
Simple Mechanical Control Systems}, Texts in Applied Mathematics, vol.~49,
Springer, New York--Heidelberg--Berlin, 2004. ISBN 0-387-22195-6.

\bibitem{SB13}
J.~W.~Simpson-Porco and F.~Bullo,
\emph{Contraction Theory on Riemannian Manifolds}, Center for Control, Dynamical
Systems, and Computation, UC Santa Barbara. Preprint submitted to
\emph{Systems \& Control Letters}, June~5, 2013. (The version cited here is the
uploaded preprint; equation and result numbers refer to it.)

\bibitem{LS98}
W.~Lohmiller and J.-J.~E.~Slotine,
\emph{On contraction analysis for nonlinear systems},
Automatica \textbf{34}(6):683--696, 1998.
[Reference~{[1]} of \cite{SB13}.]

\bibitem{LM97}
A.~D.~Lewis and R.~M.~Murray,
\emph{Configuration controllability of simple mechanical control systems},
SIAM Journal on Control and Optimization, 1997.
[Origin of the symmetric-product formalism used in \cite[Ch.~7--8]{BL04}.]

\end{thebibliography}

\end{document}